\begin{abstract}
An aggregated fleet of vehicle-to-grid (V2G) electric vehicles can
support distribution voltages and provide frequency reserve, but both
services draw on the same converter rating and stored energy. This
paper formulates their coordination as a two-time-scale control problem
in which each area aggregator chooses, every hour, a fleet power
set-point and the share of its plugged-in converter rating held as
frequency reserve. The set-points are evaluated by AC power flow on an
IEEE 33-bus feeder in each of two areas, together with an
energy-consistent fleet model with charger ratings, efficiencies,
availability, and driving demand. The reserve sets the gain and limit of
a V2G droop loop in a dynamic two-area load-frequency control model with
inertia, governors, a tie line, and automatic generation control,
simulated at a 20-ms step. A proximal policy optimization (PPO) agent is
trained with \NSeed{} independent seeds on randomized load, solar, fleet,
and disturbance scenarios and compared with uncontrolled charging, a
rule-based volt-watt and droop controller, and a day-ahead optimization
with perfect foresight, all simulated under identical conditions. On 30
held-out days, PPO lowers the daily operating cost by \TRedCostDAO\%,
\TRedCostRB\%, and \TRedCostUC\% relative to these benchmarks, reduces
undervoltage by \TRedUvPuHDAO\% and the mean frequency nadir by
\TRedNadirMeanHzDAO\% relative to the day-ahead optimization, and leaves
a smaller end-of-day energy deficit than both. The learned policy holds
reserve mainly where and when vehicles are plugged in, including
low-inertia midday hours, and releases it when converter capacity is
needed for voltage support. The energy exchanged by the frequency
droop loop is accounted for in the fleet energy balance; a comparison
with exact 3600-s simulations shows that extrapolating a 60-s window
to the hour overstates this energy by a factor of two to three, while
the treatment adopted here tracks the exact value to within 6\%, and
that no treatment changes the results. Results on out-of-distribution stress days, a reward
ablation, a weight-sensitivity study, and a bounded-action policy
comparison are also reported, and the code, trained policies, and
data are released.
\end{abstract}

\begin{keywords}
Electric vehicles, frequency regulation, load-frequency control,
proximal policy optimization, reinforcement learning, vehicle-to-grid,
voltage regulation.
\end{keywords}
